\section{Related Work}
\label{sec:related}

We situate our work within three intersecting lines of research: the evaluation of knowledge graph embedding models, statistical methods for prediction reliability, and the emerging role of reliability assessment in knowledge-based systems.

\subsection{The Knowledge Graph Evaluation Landscape}

Knowledge graph embedding models map entities and relations into continuous vector spaces~\cite{wang2017knowledge}, and large-scale benchmarking efforts~\cite{ali2021pykeen,rossi2021knowledge} have standardized evaluation protocols.
Models now span translation-based, tensor factorization, convolutional, GNN-based, hyperbolic, and LLM-enhanced paradigms~\cite{bordes2013translating,sun2019rotate,trouillon2016complex,dettmers2018convolutional,balazevic2019tucker,vashishth2020composition,zhu2021neural,le2023knowledge,lu2024deep,fang2025knowledge,gao2025unifying}.
Despite this architectural diversity, evaluation has remained anchored to MRR and Hits@$K$, introduced by~\cite{bordes2013translating}.
These metrics measure average ranking quality and have served the field well for reproducible model comparison.

\textbf{Recognized limitations.}
A growing literature has identified structural shortcomings of aggregate evaluation.
Score distributions reveal substantially more about model behavior than point estimates~\cite{reimers2017reporting}.
Systematic biases---entity degree~\cite{shomer2023toward}, hyperparameter sensitivity~\cite{lloyd2023assessing}, and benchmark artifacts~\cite{safavi2020evaluating,wang2023evaluation}---inflate apparent performance.
\cite{liu2022understanding} systematically analyzed fundamental limitations of standard protocols, and comprehensive surveys~\cite{shen2022comprehensive,li2023knowledge,jarnac2025uncertainty} have called for more rigorous evaluation methodologies.
\cite{zhu2024predictive} demonstrated that KG embedding models exhibit predictive multiplicity: re-training with different random seeds produces substantially different rankings, directly motivating statistical methods that quantify variability.
These critiques converge on a shared insight---a scalar cannot capture the full picture of model quality---but stop short of providing a systematic framework for per-prediction reliability assessment.

\textbf{Calibration and confidence.}
KG embedding scores are not calibrated probabilities.
\cite{tabacof2020probability} demonstrated this calibration failure and proposed Platt scaling as a post-hoc correction; \cite{akrami2020calibrating} independently addressed the same issue through isotonic regression.
These methods improve score interpretability but provide no formal error rate control.
\cite{jarnac2025uncertainty} surveyed the broader uncertainty landscape in KG construction and identified calibrated confidence as a critical open gap.
Our framework addresses this gap by constructing valid $p$-values that directly enable FDR control, producing per-prediction reliability guarantees at controlled error rates.

\subsection{Statistical Methods for Prediction Reliability}

\textbf{False discovery rate control.}
The FDR framework~\cite{benjamini1995controlling} revolutionized statistical inference by controlling the expected proportion of false positives among all discoveries, rather than the probability of any single false positive.
Its extensions---Storey's $q$-value with $\hat{\pi}_0$ estimation~\cite{storey2002direct,storey2003statistical} and Efron's empirical Bayes local FDR~\cite{efron2004large,efron2007size}---provide complementary perspectives on the same inferential problem.
The BH procedure guarantees FDR control under independence or the weaker PRDS condition~\cite{benjamini2001control}; a formal verification of PRDS for general KG embedding architectures is deferred to future work, though we provide empirical evidence from both the Benjamini-Yekutieli procedure (which controls FDR under arbitrary dependence) and the two-stage adaptive BH procedure~\cite{benjamini2006adaptive} in a sensitivity analysis (\S\ref{subsec:sensitivity}).
FDR methods have recently been adapted for machine learning applications, including statistical classifier comparisons~\cite{demsar2006statistical}, selective inference~\cite{taylor2015statistical,lee2016exact}, and conformalized FDR approaches~\cite{angelopoulos2023conformal,angelopoulos2024prediction,bates2024cross}, but their systematic application to KG completion evaluation has not been previously explored.

\textbf{Conformal prediction for knowledge graphs.}
A closely related line of work applies conformal prediction to KG link prediction.
\cite{marandon2024conformal} combined conformal prediction with link prediction to achieve FDR control; \cite{zhu2025conformal} applied conformal prediction to KG embeddings for prediction sets with coverage guarantees; and \cite{zhu2025certainty} introduced tailored nonconformity measures for sharper intervals.
\cite{ni2025towards} coupled conformal prediction with LLM-based KG reasoning.
Conformal methods and FDR address complementary questions:
conformal prediction produces sets with guaranteed coverage probability, answering ``what is the smallest prediction set that contains the true tail with probability $1-\alpha$?'';
FDR provides per-triple significance decisions, per-relation diagnostics, and a direct estimate of the noise proportion ($\hat{\pi}_0$)---answering ``which specific predictions are statistically reliable, and what fraction of the test set carries no signal?''
These are different inferential targets with different practical applications, and we provide the first systematic FDR-based evaluation protocol for KG completion.

\textbf{Uncertainty quantification in KG embeddings.}
Several approaches produce uncertainty estimates without formal error guarantees.
\cite{he2015learning} introduced KG2E, representing entities and relations as Gaussian distributions with variance-based uncertainty.
\cite{wu2025uncertain} proposed semi-supervised confidence distribution learning for uncertain KG completion.
These methods generate per-prediction uncertainty scores but do not control the error rate of downstream decisions---precisely the limitation that FDR control addresses.

\subsection{Reliability as a First-Class Requirement for Knowledge-Based Systems}

A distinctive thread in recent KBS scholarship emphasizes that knowledge quality is not merely a modeling concern but a system-level requirement.
\cite{purohit2025vanilla} introduced a validation framework ensuring logical consistency, integrity, and accuracy for KG completion, framing knowledge quality as a system property.
\cite{shang2024learnable} and~\cite{fang2025knowledge} advanced KG completion architectures with attention to the reliability of learned representations.
\cite{lu2024deep} developed hyperbolic embeddings and \cite{le2023knowledge} surveyed LLM-enhanced KG approaches, both acknowledging evaluation reliability as an open concern.
\cite{zhang2023conflict} tackled conflict resolution in KGs, directly engaging with the problem of distinguishing reliable from unreliable information.

Our work contributes to this trajectory by providing statistical infrastructure for knowledge reliability.
Rather than proposing a new embedding architecture or a domain-specific quality heuristic, we establish a general-purpose protocol that any KG-based system can adopt to quantify the reliability of its link predictions.
The $\hat{\pi}_0$ estimate serves as a system-level diagnostic: what fraction of this model's output should be treated as noise rather than knowledge?
This directly informs deployment decisions for knowledge-based systems in safety-critical settings.
